\newpage
\section{文献、输入来源与可复现附件}

\subsection{文献核查说明}

以下优先使用原论文、作者机构、原出版社与官方软件文档。检索截止日为 2026 年 10 月 29 日。文献支持所标明的既有事实；本报告的条件推导和反例另给证明，不把标准理论或未完成实验包装为原创定理。公开信息检索不等于穷尽全部文献，正式投稿前仍需针对最终具体贡献更新最接近方法的对照。

\begin{enumerate}[label={[R\arabic*]},leftmargin=3.1em,itemsep=7pt]
\item J. M. Lee, \emph{Introduction to Smooth Manifolds}, 2nd ed., Springer, 2013, Chapter 4. 常秩、浸入、submersion 和局部逆。\href{https://link.springer.com/chapter/10.1007/978-1-4419-9982-5_4}{出版社章节}。
\item H. J. Sussmann, ``Orbits of Families of Vector Fields and Integrability of Distributions,'' \emph{Transactions of the AMS}, 180, 171--188, 1973. \href{https://doi.org/10.1090/S0002-9947-1973-0321133-2}{DOI: 10.1090/S0002-9947-1973-0321133-2}。
\item A. S. Bandeira et al., ``Estimation under group actions: recovering orbits from invariants,'' \emph{Applied and Computational Harmonic Analysis}, 66, 236--319, 2023; preprint 2017. \href{https://arxiv.org/abs/1712.10163}{arXiv:1712.10163}。已指定群的 orbit recovery，与发现 Maxwell 对称不同。
\item G. W. Schwarz, ``Smooth functions invariant under the action of a compact Lie group,'' \emph{Topology}, 14, 63--68, 1975. \href{https://doi.org/10.1016/0040-9383(75)90036-1}{DOI: 10.1016/0040-9383(75)90036-1}。
\item S. Holman and V. Torega, ``Boundary Recovery of Anisotropic Electromagnetic Parameters for the Time Harmonic Maxwell's Equations,'' preprint 2023. \href{https://arxiv.org/abs/2303.06688}{arXiv:2303.06688}；\href{https://doi.org/10.1007/s12220-023-01499-0}{期刊 DOI}。支持各向异性参数的微分同胚 gauge 背景。
\item I. Androulidakis and G. Skandalis, ``The holonomy groupoid of a singular foliation,'' preprint 2006. \href{https://arxiv.org/abs/math/0612370}{arXiv:math/0612370}。奇异 foliation 的 groupoid 有明确向量场模条件。
\item A. Hatcher, \emph{Algebraic Topology}, Section 1.3. \href{https://pi.math.cornell.edu/~hatcher/AT/AT.pdf}{作者公开全文}。覆盖、路径提升与 monodromy。
\item P. Breiding and S. Timme, ``HomotopyContinuation.jl: A Package for Homotopy Continuation in Julia,'' \href{https://arxiv.org/abs/1711.10911}{arXiv:1711.10911}；\href{https://www.juliahomotopycontinuation.org/HomotopyContinuation.jl/stable/monodromy/}{官方 monodromy 文档}。文档包含群作用复用及数值验证限制。
\item J. D. Hauenstein and F. Sottile, ``Algorithm 921: alphaCertified: Certifying Solutions to Polynomial Systems,'' \emph{ACM Transactions on Mathematical Software}, 2012. \href{https://doi.org/10.1145/2331130.2331136}{DOI: 10.1145/2331130.2331136}。单根认证不等于根集完整性。
\item P. M. van den Berg and R. E. Kleinman, ``A contrast source inversion method,'' \emph{Inverse Problems}, 13, 1607--1620, 1997. \href{https://doi.org/10.1088/0266-5611/13/6/013}{DOI: 10.1088/0266-5611/13/6/013}。
\item X. Chen, Y. Zhong and K. Agarwal, ``Subspace methods for solving electromagnetic inverse scattering problems,'' \emph{Methods and Applications of Analysis}, 17(4), 407--432, 2010. \href{https://archive.intlpress.com/site/pub/files/_fulltext/journals/maa/2010/0017/0004/MAA-2010-0017-0004-a006.pdf}{原出版社全文}。确定/模糊电流、SOM、TSOM 和相关 Fourier 加速。
\item M. E. Everett, ``Homotopy, Polynomial Equations, Gross Boundary Data, and Small Helmholtz Systems,'' \emph{Journal of Computational Physics}, 124(2), 431--441, 1996. \href{https://doi.org/10.1006/jcph.1996.0070}{DOI: 10.1006/jcph.1996.0070}。小型离散 Helmholtz 的多根枚举先例。
\item D. W. Vasco, ``An algebraic formulation of geophysical inverse problems,'' \emph{Geophysical Journal International}, 142(3), 970--990, 2000. \href{https://doi.org/10.1046/j.1365-246x.2000.00188.x}{DOI: 10.1046/j.1365-246x.2000.00188.x}。
\item P. Li and J. Wang, ``Nonradiating sources of Maxwell's equations,'' 2024. \href{https://arxiv.org/abs/2402.10407}{arXiv:2402.10407}。源的非唯一性不直接给共同材料可实现性。
\item P. Diaconis, S. Holmes and M. Shahshahani, ``Sampling From A Manifold,'' 2012. \href{https://arxiv.org/abs/1206.6913}{arXiv:1206.6913}。coarea 条件密度与流形采样。
\item J. T. Chang and D. Pollard, ``Conditioning as disintegration,'' \emph{Statistica Neerlandica}, 51(3), 287--317, 1997. \href{https://doi.org/10.1111/1467-9574.00056}{DOI: 10.1111/1467-9574.00056}。
\item T. Cui et al., ``Likelihood-informed dimension reduction for nonlinear inverse problems,'' 2014. \href{https://arxiv.org/abs/1403.4680}{arXiv:1403.4680}。统计重要方向与精确观测商的区别。
\item O. Zahm et al., ``Certified dimension reduction in nonlinear Bayesian inverse problems,'' \emph{Mathematics of Computation}, 91, 1789--1835, 2022.\newline \href{https://arxiv.org/abs/1807.03712}{arXiv:1807.03712}。条件 likelihood profile 与认证误差条件。
\item A. Spantini et al., ``Goal-oriented optimal approximations of Bayesian linear inverse problems,'' 2017. \href{https://arxiv.org/abs/1607.01881}{arXiv:1607.01881}。任务降维的线性 Gaussian 基线。
\item L. G. Valiant, ``The Complexity of Computing the Permanent,'' \emph{Theoretical Computer Science}, 8, 189--201, 1979. \href{https://doi.org/10.1016/0304-3975(79)90044-6}{DOI: 10.1016/0304-3975(79)90044-6}。
\item M. Jerrum, A. Sinclair and E. Vigoda, ``A polynomial-time approximation algorithm for the permanent of a matrix with nonnegative entries,'' \emph{JACM}, 51(4), 671--697, 2004. \href{https://www.inf.ed.ac.uk/publications/report/0460.html}{作者机构报告}。精确困难不排除有误差保证的随机近似。
\item M. E. Everett, ``Multiple solutions to geophysical inverse problems via tracking flights across a deeply incised model landscape,'' \emph{Geophysical Journal International}, 247(3), ggag432, online 2026-10-09. \href{https://doi.org/10.1093/gji/ggag432}{DOI: 10.1093/gji/ggag432}。近期直接相关的电磁扩散多项式路径搜索先例；本文不采纳一般根集自动为循环群的推断。
\end{enumerate}

\subsection{既有科研材料的读取范围}

\textbf{[I1] Research P3L-R6。} 公共仓库 \texttt{migodam/b1-probabilistic-maxwell-imaging}，科学版本前缀 \texttt{d1bdc8aceae4c}。实际读取主报告、理论与反例、分层概率和 Track A copula 报告。\href{https://raw.githubusercontent.com/migodam/b1-probabilistic-maxwell-imaging/d1bdc8aceae4cd6a843f873b61b019d5eee3cd1d/research/p3l_r6_task_optimal_maxwell/RESEARCH_REPORT_ZH.md}{冻结主报告}；完整文件名和路径见附件 manifest。

\textbf{[I2] Research P3L-R7。} 同仓库，读取快照前缀 \texttt{975ebe7f17c2}；这标识本次读取快照，不声称它是 R7 最初发布的 freeze commit。实际读取 START\_HERE、主报告、E1 gauge audit、E2 材料纤维结果、理论、资源账本与 Go/No-Go。\href{https://raw.githubusercontent.com/migodam/b1-probabilistic-maxwell-imaging/975ebe7f17c2ba837a8cac0a5f8eb4ccf7d4e19b/research/p3l_r7_material_solution_fibers/RESEARCH_REPORT_ZH.md}{读取快照主报告}。

\textbf{[I3] 原 Research Z。} 包名为\newline
\path{Research_Z_Physics_Context_Latent_Generative_Maxwell_Imaging.zip}，2026-10-15 理论设计包。本次读取主报告、理论说明、claim ledger 与 README；完整实材料 576 维在主报告 3.A，$8\times32$ 采集在第 9 节。此包没有新的训练或 posterior 实验。

\textbf{[I4] TSOM 及其他先前结论。} SOM/TSOM 的原作者论文 [R11] 已核查。PairTSOM、F/F3 的特定项目实验在本报告中作为用户明确提供、并与已读后续报告相容的研究约束；不宣称已独立复现或完整读到其全部原始实验包。R7-T/R8-T 相关理论 PDF 已检查引用段落，其理论讨论不作为 R7 实验新增证据。

附件的 \texttt{SOURCE\_MANIFEST.json} 保存 R6/R7 完整 commit、来源路径与字节数；最终附件 manifest 再记录本次全部交付文件的 SHA256。原历史大型实验在本轮未重跑。

\subsection{可复现附件内容}

研究包包括：本报告的 LaTeX 源码；可直接执行的中文 Codex 协议；解析/monodromy toy 脚本及结果；真实二单元 TM/VIE 脚本及完整 JSON；带噪 E4 双积分脚本及结果；来源与校验 manifest；重建说明。

字体文件来自 Noto CJK 官方仓库；重建说明给出下载地址和 XeLaTeX 命令。附件不捆绑既有项目仓库或训练数据，也不修改旧轮次记录。报告 PDF 是完整科研交付，源码和结果文件用于审计与后续复现。
